% Диаграмма вариантов использования.
\begin{tikzpicture}[
  font=\footnotesize,
  uc/.style={draw, ellipse, align=center, inner xsep=3pt, inner ysep=1pt, minimum height=10mm},
  link/.style={semithick},
  dep/.style={-{Stealth[length=2mm]}, semithick, dashed},
  lbl/.style={font=\scriptsize, align=center, inner sep=1pt, fill=white},
  actor/.pic={
    \draw[semithick] (0, 0.62) circle (0.17);
    \draw[semithick] (0, 0.45) -- (0, -0.05);
    \draw[semithick] (-0.3, 0.28) -- (0.3, 0.28);
    \draw[semithick] (0, -0.05) -- (-0.22, -0.45);
    \draw[semithick] (0, -0.05) -- (0.22, -0.45);
  },
]

% --- граница системы ---
\draw[thick] (1.5, 1.0) rectangle (13.0, -10.0);
\node[anchor=north east] at (12.9, 0.95) {Игровой ассистент};

% --- варианты использования игрока ---
\node[uc] (give) at (4.5, 0) {Получить предмет\\или призвать\\существо};
\node[uc] (clarify) at (4.5, -1.95) {Ответить на\\уточняющий\\вопрос};
\node[uc] (plan) at (4.5, -3.9) {Получить план\\достижения цели};
\node[uc] (explain) at (4.5, -5.6) {Получить\\объяснение\\решения};
\node[uc] (quest) at (4.5, -7.45) {Получить задание\\с проверкой\\выполнимости};
\node[uc] (mode) at (4.5, -9.15) {Сменить режим\\работы};

% --- варианты использования администраторов ---
\node[uc] (rebuild) at (10.2, -0.75) {Пересобрать\\базу знаний};
\node[uc] (rules) at (10.2, -2.45) {Задать запреты\\и разговорные\\названия};
\node[uc] (model) at (10.2, -6.1) {Выбрать\\языковую модель};
\node[uc] (eval) at (10.2, -7.6) {Провести\\оценку качества};

% --- действующие лица ---
\pic at (0.45, -4.6) {actor};
\node at (0.45, -5.45) {Игрок};
\coordinate (p) at (0.7, -4.35);

\pic at (14.8, -1.55) {actor};
\node[align=center] at (14.8, -2.55) {Администратор\\базы знаний};
\coordinate (a1) at (14.5, -1.3);

\pic at (14.8, -6.75) {actor};
\node[align=center] at (14.8, -7.75) {Администратор\\системы};
\coordinate (a2) at (14.5, -6.5);

% --- связи игрока ---
\foreach \target in {give, plan, explain, quest, mode}
  \draw[link] (p) -- (\target.west);

% --- связи администраторов ---
\draw[link] (a1) -- (rebuild.east);
\draw[link] (a1) -- (rules.east);
\draw[link] (a2) -- (model.east);
\draw[link] (a2) -- (eval.east);

% --- отношения между вариантами ---
\draw[dep] (clarify) -- node[lbl, right=2pt] {<<расширяет>>} (give);
\draw[dep] (plan) -- node[lbl, right=2pt] {<<включает>>} (explain);

\end{tikzpicture}
